%% ma: L11-B14-0021
%% lop: 11
%% chuong: 4
%% bai: 14
%% dang: 11.14.2
%% loai: DS
%% muc: [NB, NB, TH, NB]
%% dap_an: "SĐĐS"
%% nguon: "tai-lieu-hoc-tap-mon-toan-lop-11-theo-chuong-trinh-giao-duc-pho-thong-2018.pdf (Lớp Toán Thầy Duy, Phạm Lê Duy, THPT Chu Văn An), Chương 4 Bài 5, Phần 2, Câu 1 (tr.532)"
%% kien_thuc: "Lớp 11 Bài 14 (phép chiếu song song)"
%% trang_thai: chinh-thuc
%% ghi_chu: "Đã duyệt 03/11/2026."
\begin{ex}
Cho các mệnh đề sau?
\choiceTF
{Hình chiếu song song của hai đường thẳng song song là hai đường thẳng cắt nhau.}
{\True Phép chiếu song song biến ba điểm thẳng hàng thành ba điểm thẳng hàng và không làm thay đổi thứ tự giữa các điểm.}
{\True Cho mặt phẳng $(P)$, đường thẳng $l$ cắt $(P)$ và điểm $M$ không nằm trên đường thẳng $l$. Gọi điểm $M^{\prime}$ là ảnh của điểm $M$ lên mặt phẳng $(P)$ theo phương chiếu $l$ của phép chiếu song song. Khi đó $MM^{\prime}$ song song $l$.}
{Phép chiếu song song làm thay đổi tỉ số độ dài của hai đoạn thẳng nằm trên hai đường thẳng song song hoặc trùng nhau.}
\loigiai{\sloppy a) Hình chiếu của hai đường thẳng song song là hai đường thẳng song song hoặc trùng nhau. Sai.
b) Đúng (tính chất 1).
c) $M^{\prime}$ là giao của $(P)$ với đường thẳng qua $M$ song song $l$, nên $MM^{\prime}\parallel l$. Đúng.
d) Phép chiếu song song không làm thay đổi tỉ số độ dài. Sai.}
\end{ex}
